\abstract{Digital serious games for children with attention-deficit/hyperactivity disorder (ADHD) can improve selected trained attention and executive-function outcomes, but clinically meaningful transfer beyond the trained task remains a central challenge. ARC-ADHD is proposed as a closed-loop serious-game framework in which controlled executive challenge is followed by an explicit pause, orientation to current bodily or task state, strategy selection, regulation, recovery monitoring, and return to the interrupted goal. The core intervention deliberately excludes augmented-reality transfer experiments, facial-behavior feedback, bilateral auditory stimulation, and wearable EEG acquisition, which are assigned to separate companion studies or measurement tracks. The principal mechanistic quantity is monitoring-anchored recovery latency \(T_{\mathrm{rec}}\). Recovery is operationalized against a stable target-rule calibration reference using a small preregistered behavioral/timing feature space and a persistence criterion; the immediate preswitch window is used for challenge eligibility rather than as the recovery normalization baseline. Phase 0 first tests recovery-measure validity through prespecified computability, scale, timing-coverage, and perturbation-stability gates before the selected definition is used inferentially. The paper then distinguishes proximal mechanism change, independent near transfer, and externally assessed far transfer, with candidate-specific event denominators and explicit right-censoring. The staged programme comprises co-design and construct validation, feasibility testing, a pilot randomized trial intended to estimate design parameters rather than establish efficacy, and a later confirmatory trial. EDI Reactivity is the leading provisional external endpoint candidate for Phase 1/2 planning, while the confirmatory endpoint remains unfrozen pending a defensible Polish localization and validation pathway. ARC-ADHD is presented as a falsifiable and reproducible computational research framework rather than an established treatment or a completed randomized-trial protocol.}
